\documentclass{article}
\usepackage{amsmath,amssymb,amsthm,mathtools}
\newtheorem{theorem}{Theorem}
\newtheorem{lemma}{Lemma}
\newtheorem{proposition}{Proposition}
\newtheorem{corollary}{Corollary}
\newtheorem{definition}{Definition}
\newtheorem{example}{Example}
\title{ShadowBench geometry/L2/geo\_gen\_L2\_005}
\begin{document}
\maketitle
% Problem id: geometry/L2/geo_gen_L2_005
% Companion instructions: docs/instructions.md
% Candidate skeletons: docs/skeletons/
% Lean target: ShadowBench/Source/Main.lean

\begin{theorem}[monges_circle_theorem]\label{thm:monge}
Monge's theorem states that for any three circles in a plane, none of which is completely inside one of the others, the intersection points of each of the three pairs of external tangent lines are collinear.

In this formalization, we consider the affine-plane case where the three radii are pairwise distinct, so that these intersection points are finite points rather than points at infinity. Degenerate circles of radius zero are allowed, provided the radii remain pairwise distinct.

For any two circles in a plane, an external tangent is a line that is tangent to both circles but does not pass between them. There are two such external tangent lines for any two circles satisfying the above non-containment assumptions. Monge's theorem states that the three such points given by the three pairs of circles lie on a straight line.
\end{theorem}
\end{document}
